\subsection*{2.2}\phantomsection \label{sga3:VIIB:2.2}\addcontentsline{toc}{subsection}{\SGARefLink{sga3:VIIA:2.2}{2.2}. \(\mathbf O_k\) 上的双代数}

设 \(\mathbf H\) 为 \(\mathbf O_k\) 上余代数范畴中的群对象。于是对 \(\operatorname{Alf}/k\) 的每个对象 \(C\)，对象 \(\mathbf H(C)\) 具有群中的 \(C\)-余代数结构 \textup{（参见 VII\textsubscript{A} \SGARefLink{sga3:VIIA:3.2}{3.2}）}。遵循 Manin，我们用\emph{双代数}\footnote{\label{sga3:VIIB:nde:77}也使用术语“bigebra” \textup{（参见 [BAlg]，III \S~11.4）}。回顾 \textup{（参见 VII\textsubscript{A}，\SGARefLink{sga3:VIIB:3.1:lemma}{3.1}，编者注~(26)）}，这里考虑的所有双代数均假设为余交换且带有对极，因而是余交换 Hopf 代数。}来代替“群中的余代数”。此外，对 \(\operatorname{Alf}/k\) 中的每个态射 \[
  \phi:C\longrightarrow D
\]，映射 \[
  D\otimes_C\mathbf H(C)\longrightarrow\mathbf H(D)
\] 是 \(D\)-双代数同态。

\medskip\noindent\textbf{定义。}\phantomsection \label{sga3:VIIB:2.2.definition}将这些性质概括地表述为：\(\mathbf H\) 是\emph{\(\mathbf O_k\) 上的双代数}。

\SGARefLink{sga3:VIIB:1.3.5}{1.3.5} 的函子 \[
  \mathbf H\longmapsto\operatorname{Spf}^*(\mathbf H)
\] 显然与有限乘积交换。因此它将 \(\mathbf O_k\) 上的双代数映为 \(k\)-群函子，即从 \(\operatorname{Alf}/k\) 到群范畴的协变函子。

事实上，对每个有限长度 \(k\)-代数 \(C\)，\[
\begin{aligned}
  \operatorname{Spf}^*(\mathbf H)(C)
  &=
  \operatorname{Spf}^*(\mathbf H(C))
  \\
  &=
  \bigl\{
    x\in\mathbf H(C)
    \mid
    \varepsilon(x)=1
    \ \text{且}\
    \Delta(x)=x\otimes x
  \bigr\}
\end{aligned}
\] 的元素在代数 \(\mathbf H(C)\) 的乘法下构成群；参见 VII\textsubscript{A} \SGARefLink{sga3:VIIA:3.2.2}{3.2.2}。还需注意，\(\Delta(x)=x\otimes x\) 蕴含 \[
  \varepsilon(x)=\varepsilon(x)^2.
\]。所以当 \(C\) 为局部环且 \(x\neq0\) 时，有 \(\varepsilon(x)=1\)。\footnote{\label{sga3:VIIB:nde:78}由于 \(\Delta(x)=x\otimes x\)，有 \(x=\varepsilon(x)x\)；故只有当 \(x=0\) 时才可能有 \(\varepsilon(x)=0\)。}
